% GENERATED FILE. Edit canonical lessons, then run npm run generate:books.
% canonical-source-hash: fnv1a64:11070c41d757a1f5
% canonical-lessons: IT-C472-testardo, IT-C472-pigro, IT-C472-generoso, IT-C472-avaro, IT-C472-vantarsi

\chapter{Pregi e difetti}
\label{ch:it-pregi-e-difetti}
\hlchaptermodality{472}

\begin{chapteropening}
\textbf{By the end of this chapter:} \emph{I can describe the good and bad sides of someone's character.}

Complete the last lesson of chapter 472: I can describe the good and bad sides of someone's character.
\end{chapteropening}

\section[testardo]{testardo — stubborn}

\label{lesson:IT-C472-testardo}



\begin{quote}
Before the new one: say the Italian for chatterbox, then the Italian for to blush.
\end{quote}

\subsection*{You'll want to know: testardo}
\textbf{testardo} — \textquotedblleft{}stubborn\textquotedblright{}.

\begin{quote}
Mio nonno è testardo: non cambia mai idea.
\end{quote}

\begin{quote}
\emph{My grandfather is stubborn: he never changes his mind.}
\end{quote}

From \emph{testa} \textquotedblleft{}head\textquotedblright{}: someone with a hard head. \emph{Avere la testa dura} means the same.

\begin{cousinweb}
\emph{testardo} is built on \emph{testa} \textquotedblleft{}head\textquotedblright{}, which is inherited from Latin \emph{testa} \textquotedblleft{}an earthen pot, a shell\textquotedblright{} (EtymDB; Lewis \& Short); its later sense \textquotedblleft{}skull\textquotedblright{} became the everyday word for head.
\end{cousinweb}

\begin{grammarlens}[title={What you've built}]
One more word for telling, explaining and describing at B1.
\end{grammarlens}

\subsection*{Your turn}
\begin{itemize}
\raggedright
  \item \emph{Say it:} \emph{testardo}
  \item \emph{Say it:} \emph{testardo}, once more
  \item \emph{Say it:} a sentence of your own with \emph{testardo}
  \item \emph{Recall:} say the Italian for chatterbox, then the Italian for to blush, then say \emph{testardo} again
\end{itemize}

\begin{tcolorbox}[breakable,colback=teal!4,colframe=teal!35!black,title={Before you move on}]
What does testardo mean? (\textquotedblleft{}Stubborn\textquotedblright{}.) Say it once more.
\end{tcolorbox}
\section[pigro]{pigro — lazy}

\label{lesson:IT-C472-pigro}



\begin{quote}
Before the new one: say the Italian for to blush, then the Italian for stubborn.
\end{quote}

\subsection*{You'll want to know: pigro}
\textbf{pigro} — \textquotedblleft{}lazy\textquotedblright{}.

\begin{quote}
La domenica sono pigro e resto a letto fino a tardi.
\end{quote}

\begin{quote}
\emph{On Sundays I am lazy and stay in bed until late.}
\end{quote}

The noun is \emph{la pigrizia} \textquotedblleft{}laziness\textquotedblright{}.

\begin{cousinweb}
\emph{pigro} is inherited from Latin \emph{piger} \textquotedblleft{}reluctant, slow\textquotedblright{} (EtymDB), related to \emph{pigēre} \textquotedblleft{}to feel reluctant\textquotedblright{} (Lewis \& Short).
\end{cousinweb}

\begin{grammarlens}[title={What you've built}]
One more word for telling, explaining and describing at B1.
\end{grammarlens}

\subsection*{Your turn}
\begin{itemize}
\raggedright
  \item \emph{Say it:} \emph{pigro}
  \item \emph{Say it:} \emph{pigro}, once more
  \item \emph{Say it:} a sentence of your own with \emph{pigro}
  \item \emph{Recall:} say the Italian for to blush, then the Italian for stubborn, then say \emph{pigro} again
\end{itemize}

\begin{tcolorbox}[breakable,colback=teal!4,colframe=teal!35!black,title={Before you move on}]
What does pigro mean? (\textquotedblleft{}Lazy\textquotedblright{}.) Say it once more.
\end{tcolorbox}
\section[generoso]{generoso — generous}

\label{lesson:IT-C472-generoso}



\begin{quote}
Before the new one: say the Italian for stubborn, then the Italian for lazy.
\end{quote}

\subsection*{You'll want to know: generoso}
\textbf{generoso} — \textquotedblleft{}generous\textquotedblright{}.

\begin{quote}
È una persona generosa e aiuta sempre gli amici.
\end{quote}

\begin{quote}
\emph{She is a generous person and always helps her friends.}
\end{quote}

Of food: \emph{una porzione generosa} — a big helping.

\begin{cousinweb}
\emph{generoso} is a learned borrowing from Latin \emph{generōsus} \textquotedblleft{}of noble birth, noble-minded\textquotedblright{} (EtymDB; Lewis \& Short), built on \emph{genus} \textquotedblleft{}birth, family\textquotedblright{}.
\end{cousinweb}

\begin{grammarlens}[title={What you've built}]
One more word for telling, explaining and describing at B1.
\end{grammarlens}

\subsection*{Your turn}
\begin{itemize}
\raggedright
  \item \emph{Say it:} \emph{generoso}
  \item \emph{Say it:} \emph{generoso}, once more
  \item \emph{Say it:} a sentence of your own with \emph{generoso}
  \item \emph{Recall:} say the Italian for stubborn, then the Italian for lazy, then say \emph{generoso} again
\end{itemize}

\begin{tcolorbox}[breakable,colback=teal!4,colframe=teal!35!black,title={Before you move on}]
What does generoso mean? (\textquotedblleft{}Generous\textquotedblright{}.) Say it once more.
\end{tcolorbox}
\section[avaro]{avaro — mean, stingy; miser}

\label{lesson:IT-C472-avaro}



\begin{quote}
Before the new one: say the Italian for lazy, then the Italian for generous.
\end{quote}

\subsection*{You'll want to know: avaro}
\textbf{avaro} — \textquotedblleft{}mean, stingy; miser\textquotedblright{}.

\begin{quote}
Non è avaro, ma non spende per cose inutili.
\end{quote}

\begin{quote}
\emph{He is not stingy, but he does not spend on useless things.}
\end{quote}

Also a noun: \emph{un avaro} — a miser.

\begin{cousinweb}
\emph{avaro} is inherited from Latin \emph{avārus} \textquotedblleft{}greedy\textquotedblright{} (EtymDB), from \emph{avēre} \textquotedblleft{}to long for, to crave\textquotedblright{} (Lewis \& Short).
\end{cousinweb}

\begin{grammarlens}[title={What you've built}]
One more word for telling, explaining and describing at B1.
\end{grammarlens}

\subsection*{Your turn}
\begin{itemize}
\raggedright
  \item \emph{Say it:} \emph{avaro}
  \item \emph{Say it:} \emph{avaro}, once more
  \item \emph{Say it:} a sentence of your own with \emph{avaro}
  \item \emph{Recall:} say the Italian for lazy, then the Italian for generous, then say \emph{avaro} again
\end{itemize}

\begin{tcolorbox}[breakable,colback=teal!4,colframe=teal!35!black,title={Before you move on}]
What does avaro mean? (\textquotedblleft{}Mean, stingy; miser\textquotedblright{}.) Say it once more.
\end{tcolorbox}
\section[vantarsi]{vantarsi — to boast, to show off}

\label{lesson:IT-C472-vantarsi}



\begin{quote}
Before the new one: say the Italian for generous, then the Italian for mean.
\end{quote}

\subsection*{You'll want to know: vantarsi}
\textbf{vantarsi} — \textquotedblleft{}to boast, to show off\textquotedblright{}.

\begin{quote}
Il mio collega si vanta sempre della sua macchina nuova.
\end{quote}

\begin{quote}
\emph{My colleague is always boasting about his new car.}
\end{quote}

\emph{Vantarsi di qualcosa}: the thing you boast about takes \emph{di}.

\begin{cousinweb}
\emph{vantarsi} is usually traced to Late Latin \emph{vānitāre} \textquotedblleft{}to speak emptily\textquotedblright{}, built on Latin \emph{vānus} \textquotedblleft{}empty, vain\textquotedblright{} (Lewis \& Short).
\end{cousinweb}

\begin{grammarlens}[title={What you've built}]
One more word for telling, explaining and describing at B1.
\end{grammarlens}

\subsection*{Your turn}
\begin{itemize}
\raggedright
  \item \emph{Say it:} \emph{vantarsi}
  \item \emph{Say it:} \emph{vantarsi}, once more
  \item \emph{Say it:} a sentence of your own with \emph{vantarsi}
  \item \emph{Recall:} say the Italian for generous, then the Italian for mean, then say \emph{vantarsi} again
\end{itemize}

\begin{tcolorbox}[breakable,colback=teal!4,colframe=teal!35!black,title={Before you move on}]
What does vantarsi mean? (\textquotedblleft{}To boast, to show off\textquotedblright{}.) Say it once more.
\end{tcolorbox}
